% !PRINTED 555
\weberpage{555}
%<<CONT
lies, the finiteness of the number of classes is thereby proved.

If, as hitherto, we let $\omega_1, \omega_2, \ldots, \omega_n$ denote a minimal basis of $\Omega$, then all the integers of the field are obtained from

\weq{1}{\omega = \omega_1 x_1 + \omega_2 x_2 + \cdots + \omega_n x_n}

when we assign integral rational numerical values to the $x_1, x_2, \ldots, x_n$. If we assume a positive integer $k$ and stipulate that none of the numbers $x_i$ shall step outside the interval $\pm k$, then the \gesperrt{absolute value} of $\omega$ will not lie above the limit

\[
(|\omega_1| + |\omega_2| + \cdots + |\omega_n|) k = rk
\]

provided that by $|\omega_1|, |\omega_2|, \ldots$ are understood the absolute values (Einleitung, Bd. I, S. 18) of the (real or complex) quantities $\omega_i$, and $r$ denotes the sum $|\omega_1| + |\omega_2| + \cdots$. If we form the expression (1) for the $n$ conjugate fields, and take the product, then, denoting by $R$ a positive real number that lies above the product of the $n$ values $r$, we obtain

\weq{2}{N_a(\omega) < R k^n.}

This number $R$ depends only on the nature of the field $\Omega$, and not on $k$.

Now let $\mu$ be any integral functional in $\Omega$, and $N_a(\mu)$ its absolute norm. If we determine the integer $k$ so that

\weq{3}{k^n \leqq N_a(\mu) < (k+1)^n,}

and if we further assign to the numbers $x_i$ in (1) the values $0, 1, 2, \ldots, k$, then the number of the distinct values that arise from (1) is $(k+1)^n$, and hence greater than $N_a(\mu)$. But by \S.\,148, 3. the number of the numbers incongruent with respect to the modulus $\mu$ is equal to $N_a(\mu)$, and consequently among the numbers $\omega$ so determined there must occur at least two distinct numbers congruent with respect to the modulus $\mu$. If then $\omega' \equiv \omega'' (\mathrm{mod}\ \mu)$, the difference

\weq{4}{\alpha = \omega' - \omega'' = (x_1' - x_1'') \omega_1 + \cdots + (x_n' - x_n'') \omega_n}

will be divisible by $\mu$, and at the same time the integers

\[
x_1' - x_1'' = a_1, \ldots, x_n' - x_n'' = a_n
\]
%CONT>>
